\chapter{Normas Técnicas}

\section{Normas de Desarrollo y Documentación de Software}

Se ha seguido un enfoque estructurado de desarrollo, utilizando buenas prácticas de codificación y organización del código fuente. En particular:

\begin{itemize}
    \item PEP 8 -- Guía de estilo para Python, aplicable al código implementado para el entrenamiento y despliegue del modelo \cite{pep82025}.
\end{itemize}

\section{Normas sobre gestión de datos}

El uso del corpus OffendES, aunque público, sigue prácticas de manejo ético de datos. Por tanto, se han aplicado normas como:

\begin{itemize}
    \item ISO/IEC 27001 -- Gestión de seguridad de la información \cite{iso27001}.
\end{itemize}

\section{Normas sobre modelos de IA}

En el desarrollo del modelo de lenguaje ofensivo se ha considerado:

\begin{itemize}
    \item Guía ética para inteligencia artificial de la UNESCO (2021), enfocada en la transparencia, no discriminación y trazabilidad \cite{unesco2024ethics}.
\end{itemize}

% =========================
% Página 67
% =========================
